\chapter{Процессы роста малых кластеров металла и их теплообмена с окружающей средой}\label{ch:ch6}

Наночастицы металлов являются интересным объектом исследований с множеством практических применений. К ним относится и их использование в качестве катализаторов, и улучшенние свойств масел и смазок, и потеницальное применение для точной доставки лекарств.

Однако, свойства металлических наночастиц зависят от их размера, который в свою очередь определяется отводом тепла. На эту характеристику влияет коэффициент термической аккомодации, который показывает эффективность процесса теплопередачи от наночастиц к буферному газу. Эта глава посвящена вычислению этой величины методом молекулярной динамики.

В первом подразделе представлена исследуемая система, вывод её на равновесие и описан метод расчета коэффцицента термической аккомодации и вероятности роста кластера. Во втором представлены результаты расчета коэффициента термической аккомодации и теплоотдачи для металлических кластеров. В третьем показаны данные по вероятности прилипания атома к кластеру. В четверотом подразделе обсуждается сравнение с экспериментальными данными и иными методами расчета коэффициента термической аккомодации. Результаты, приведенные в данной главе, опубликованы в~\cite{__2019, lenev_thermal_2016, lenev_sticking_2018}.

\section{Описание исследуемой системы}

Цеетральной частью исследуемой системы является кластер, состоящий из атомов железа. Он описывается многочастичным потенциалом Финниса-Синклера~\cite{mendelev_development_2003}, который явлется вариантом потенциала EAM:

\begin{align}
	U = U_{pair}(r) + U_{manybody}(\rho_{el}(r)),
\end{align}

где $U_{pair}$ \--- энергия парных взаимодействий, зависящая от расстояния между атомами, $U_{manybody}$ \--- "многочастичная часть", которая зависит от $\rho_{el}$ \--- электронной плотности. Значения парной энергии и электронной плотности задаются таблично. В отличие от EAM, потенциал Финниса-Синклера учитывает тип атома, который вносит вклад в электронную плотность. Данный потенциал хорошо моделирует процесс плавления в металлах с объемноцентрированной кристаллической решеткой. 

Рассмотрены кластеры из 9, 27 и 113. Количество атомов в кластере определялось радиусом области, в которой размещались атомы железа в узлах кристаллической решетки. Начальные скорости атомов задаются согласно распределению Максвелла для 300~К. Далее кластер нагревается с помощью термостата Ланжевена до нужной температуры в течение 20~пс. Чтобы в силу стохастического характера этого термостата кластер не приобретал скорость, расчет происходит в системе отсчета, связанной с центром масс кластера. Под температурой кластера в работе подразумевается следующее выражение:

\begin{align}
	M \langle v^2 \rangle /2 = 3 R T, T = M \langle v^2 \rangle / 3 R,
\end{align}

где $M$ \--- молярная масса железа, $\langle v^2 \rangle$ \--- среднеквадратичная скорость атомов, $R$ \--- универсальная газовая постоянная.

Рассмотрены температуры кластера от 200 до 2500~К. При более высоких температурах может наблюдаться распад кластера. Также при такой температуре находится металлический пар в некоторых экспериментах по нуклеации. После достижения нужной температуры моделирование продолжается ещё 5~пс, чтобы удостовериться в наличии равновесия. 

После вывода на равновесие в систему добавляется налетающий атом. Взаимодействие между атомом и кластером описывается потенциалом Букингема:

\begin{align}
	U(r) = A e^{-r/B} - C/r^6,
\end{align}

где $A = 4617~eV$, $B=0.2752$~\AA, $70$~eV/\AA$^6$. Здесь притяжение аналогично потенциалу Леннард-Джонса, а для отталкивания используется более сильная зависимость. Радиус обрезания, когда потенциал взаимодействия обнуляется $r_{cut} = 2.5~\sigma = 7.165$~\AA. Здесь под $\sigma$ имеется в виду расстояние, на котором потенциал проходит через ноль.

Начальное положение атома располагается на дуге как показано на рисунке~\cref{fig:model_accommodation}. Варьируется прицельный параметр $r$ таким образом, чтобы получить 20-25 точек в области, где имеет место взаимодействие с кластером. Такое количество точек позволяет как получить приемлемую погрешность для каждой точки, так и итоговая погрешность была меньше 0.3 от величины. Ось $x$ выбрана таким образом, чтобы все начальные положения атома имели координату $x = 0$. Ось $y$ проведена так, чтобы $y(t=0) = r$. На рисунке $R$ \--- радиус области, в которую изначально помещаются атомы кластера, служит характерным изображением кластера. Начальная скорость $v_i$ направлена по оси $z$ и выбирается случайным образом так, чтобы использованный набор скоростей соответствовал распределению Максвелла для 300~К. 

\begin{figure}[ht]
    \centerfloat{
        \includegraphics[scale=0.8]{model.png}
    \caption{Пример схемы численного эксперимента для кластера, состоящего из 27 атомов. Окружность радиуса R - упрощённое изображение центров внешних атомов кластера. Кружок – налетающий атом: r – прицельный параметр, $v_0$ – начальная скорость, тонкая линия со стрелкой - траектория налетающего атома. Пунктирная дуга – геометрическое место начальных положений налетающего атома.}
    }\label{fig:model_accommodation}
\end{figure}

На рисунках~\cref{fig:trajectories} показаны две характерные траектории атомов. На первом рисунке~\cref{fig:traj1} показана траектория с большим количеством соударений налетающего атома с кластером. Это видно по количеству двойных пиков, каждый из которых соответствует одному соударению. Вначале атом, приближаясь к кластеру, проходит потенциальную яму и набирает скорость(первый пик), затем приближется максимально близко, где отражается от поверхности кластера, имея минимальную скорость, и затем удаляется(второй пик). Поскольку кластер и атом имеют близкую и низкую относительно температуры кипения скорость, атому требуется большое время и много соударений, чтобы набрать энергию, необходимую для отрыва от кластера. Во втором случае~\cref{fig:traj2} кластер имеет высокую температуру, и одного соударения хватает для достижения необходимой энергии. 


\begin{figure}[ht]
    \centerfloat{
        \hfill
        \subcaptionbox[List-of-Figures entry]{\label{fig:traj1}}{%
            \includegraphics[width=0.5\linewidth]{accommodation_traj1.png}}
        \hfill
        \subcaptionbox{\label{fig:traj2}}{%
            \includegraphics[width=0.5\linewidth]{accommodation_traj2.png}}
    }
    \caption{Два примера молекулярно-динамических траекторий для взаимодействия налетающего атома с кластером из
9 атомов при температурах кластера 500 К (а) и 2500 К (б).}\label{fig:trajectories}
\end{figure}

В начале и конце зависимости модуля скорости от времени имеется плато. Оно соответствует области, в которой атом и кластер не взаимодействуют, и по нему определяются конечная и начальная скорость атома, из которых рассчитывается, какую энергию приобрел атом за время взаимодействия с кластером. По рисункам видно, что в случае одного соударения с "горячим" кластером эта величина больше, однако, далее будет показано, чо за счет краткости взаимодействия эффективность передачи тепла ниже.

\section{Моделирование термической аккоммодации и оценка теплоотдачи}

Рассматривается процесс передачи энергии атомам аргона кластеру железа. Первой из величин, представляющих интерес, является количество полученной атомом энергии. Эта величина показывает, насколько нагрелся или остыл кластер в результате взаимодействия. Эта величина рассчитывается для каждой траектории и усредняется для каждого значения прицельного параметра:

\begin{align}
	Q = m \langle v_f^2 - v_i^2 \rangle_r /2k_B,
\end{align}

где $m$ \--- атомная масса, $v_f$ и $v_i$ \--- начальная и конечная скорости налетающего атома соответственно, $k_B$ \--- константа Больцмана. Усреднение проводится по набору траекторий для каждого значения прицельного параметра. Погрешности, показанные на рисунке, рассчитанны как корень из дисперсии для данной серии значений. 
Всего для каждой температуры для кластера размером 9 атомов было рассмотрено 300 000 траекторий, для кластера из 27 атомов \--- 100 000, для кластера из 113 атомов \--- 10 000.Для удобства и наглдяности величина переданной атомом энергии дается в К, поскольку её нужно сравнивать с разностью между температурами налетающего атома и кластера. 
было рассмотрено 300 000 траекторий, для кластера из 27 атомов \--- 100 000, для кластера из 113 атомов \--- 10 000.

\begin{figure}[ht]
    \centerfloat{
        \includegraphics[scale=0.8]{aiming_parameter.png}
    \caption{Зависимость количества полученной атомом энергии от прицельного параметра для N = 27 и T = 1500~K. Горизонтальная линия показывает усредненное по плато значение, вертикальная – значение $r_{\sigma}$, стрелка – среднее расстояние от центра масс кластера до наиболее удаленного атома. Температура потока налетающих атомов 300~K.}
    }\label{fig:aiming_parameter}
\end{figure}

На зависимости величины переданной атому энергии от прицельного параметра~\cref{fig:aiming_parameter} можно увидеть несколько участков. На первом из них эта величина флуктуирует относительно некоторой средней величины. Усреднением полученных результатов на этом участке получается значение теплоты, переданной атому, для данных температуры и размера кластера. 

Далее на зависимости от прицельного парамера наблюдается спад, который объясняется ослаблением взаимодействия между налетающим атомом и кластером. Можно заметить, что спад начинается не при значениях прицельного параметра, равных характерному радиусу кластера, а более близких к сумме радиуса кластера и расстояния, на котором потенциальная энергия взаимодействия атома и кластера принимает минимальное значение.

Последний участок зависимости переданной атому энергии от прицельного параметра является нулевой, на котором отсутствует взаимодействие атома и кластера, и потому этот участок не представляет интереса.

Для усреднения при каждых размере кластера и температуре использовалась следующая формула:

\begin{align}
	Q \sigma = \int_0^{\infty} Q(r) 2\pi r dr,
\end{align}

где $\sigma$ \--- площадь сечения, определяемая как $\pi r_{\sigma}^2$, где $r_{\sigma}$ соответствует середине участка спада на графике зависимости величины полученной атомом энергии от прицельного параметра. По значениям радиуса сечения из таблицы можно сделать вывод, что он не зависит от температуры. Зависимость можно упрощенно представить в форме ступеньки, которая заканчивается на радиусе сечения взаимодействия. В таком случае высота ступеньки будет нужным средним значением теплоты, переданной атому кластером.

Зависимость осредненного количества теплоты от температуры является возрастающей~\cref{fig:Q_T}, то есть чем выше температура кластера, тем больше энергии забирает налетающий атом. В случае температуры кластера, равной 200~К, количество полученной атомом теплоты отрицательно, следовательно, атом в среднем отдает энергию, что логично, потому что поток атомов имеет большую температуру, чем кластер. Можно заметить, что температурная зависимость проходит через ноль примерно при температуре 300~К, равной температуре кластера. Следовательно, в этом случае среднее количество полученной атомом энергии равно нулю. 

\begin{figure}[ht]
    \centerfloat{
        \includegraphics[scale=0.8]{Q_T.png}
    \caption{Зависимость усредненного количества полученной атомом энергии от температуры кластера для трех размеров кластера: N = 9, 27, 113. Температура потока налетающих атомов 300~K.}
    }\label{fig:Q_T}
\end{figure}

Полученная зависимость изменения энергии атома от температуры возрастает слабее, чем линейная зависимость. Для выяснения причин этого явления была получена зависимость времени, которое атом проводит возле кластера, от температуры~\cref{fig:t_T}. Это время было получено следующим образом. Отмечался момент времени $t_1$, когда сила действующая на атом, начинала отличаться от нуля. Затем отмечался момент $t_2$, когда эта сила вновь становилась равна нулю. То есть определялось время, которое атом проводил в зоне действия потенциала кластера. Эта зависимость является убывающей. Чем выше температура кластера, тем меньше времени атом взаимодействует с кластером и меньше приобретает энергии. Таким образом, чем дольше атом проводит время в зоне действия потенциала кластера и испытывает соударений, тем ближе его конечная температура к температуре кластера и тем выше эффективность отвода тепла.

\begin{figure}[ht]
    \centerfloat{
        \includegraphics[scale=0.8]{t_T.png}
    \caption{Зависимость времени взаимодействия для N = 113 от температуры кластера. Температура потока налетающих атомов 300~K.}
    }\label{fig:t_T}
\end{figure}

Эту эффективность характеризует коэффициент термической аккомодации, определяемый формулой:

\begin{align}
	\beta = Q/2(T-T_a),
\end{align}

где $T$ и $T_a$ \--- температуры кластера и налетающего атома соответственно. Выражение представляет собой отношение фактически полученной атомом энергии к оценке максимально возможной энергии, которую может получить атом аргона. Этот коэффициент строго больше нуля при достаточном усреднении. Если атом горячее кластера, он отдает энергию, и тогда и числитель, и знаменатель меньше нуля. В противном случае оба имеют положительный знак. Также коэффициент термической аккомодации меньше или приблизительно равен 1, поскольку атом не может нагреться до температуры выше более горячего кластера или остыть сильнее более холодного кластера. Максимальное значение, 1, получается, если атом находится рядом с кластером достаточно долго, множество раз соударяется с ним и в результате получает температуру, близкую к температуре кластера.

Следует заметить, что описанные выше свойства выполняются только при усреднении достаточного количества траекторий. Для отдельных траекторий в силу случайности выбора скоростей, например, скорость налетающего атома может быть выше температуры кластера при том, что температура потока $T_a$ будет ниже. Тогда атом потеряет энергию при положительном знаменателе, коэффициент термической аккомодации будет меньше нуля. Если же скорость атома будет значительно выше средней по потоку и выше температуры кластера, то коэффициент $\beta$ окажется выше единицы.

Так как зависимость полученной от кластера атомом энергии возрастает слабее, чем линейно, коэффициент термической аккомодации убывает с ростом температуры~\cref{fig:alpha_T}. Поскольку для покидания зоны действия потенциала кластера необходимо преодолеть потенциальный барьер, было сделано предположение об активационном характере зависимости. Для этого была сделана аппроксимация и построена зависимость коэффициента термической аккомодации от обратной температуры в полулогарифмическом масштабе. Видно, что полученные результаты не противоречат этой гипотезе. Следует также заметить, что коэффициенты наклона различны для разных размеров кластеров.

\begin{figure}[ht]
    \centerfloat{
        \includegraphics[scale=2.0]{alpha_T.png}
    \caption{Зависимость коэффициента термической аккомодации для N = 9, 27, 113  от обратной температуры. Температура потока налетающих атомов 300~K. }
    }\label{fig:alpha_T}
\end{figure}

\section{Сравнение с экспериментальными результатами и расчетами коэффициента аккомодации для плоской поверхности}

Для некоторых величин, характеризующих кластеры, свойственна зависимость от размера кластера, пропорциональная $N^{-1/3}$~\cite{smirnov_processes_2013, smirnov_cluster_2013}. Первой из таких величин является $A$ \--- предположительно величина потенциального барьера. Она возрастает с ростом кластера. Причиной этого служит увеличение числа атомов, с которыми взаимодействует налетающий атом.

Порядок величины энергии активации можно независимо определить по графику зависимости скорости от времени. Поскольку максимумы кинетической энергии соответствуют минимумам потенциальной, высота последнего максимума будет равна величине потенциального барьера. Для траекториий 1 и 2 получаются значения 200~К и 300~К соответственно. Следует заметить, что эта оценка не может являться точной, поскольку проведена только по одной из траекторий, хотя и характерной.

\begin{figure}[ht]
    \centerfloat{
        \includegraphics[scale=0.8]{A_N.png}
    \caption{Зависимость энергии активации от размера кластера.}
    }\label{fig:A_N}
\end{figure}

Поскольку зависимости $ln \beta (T)$ и $A(N^{-1/3})$~\cref{fig:A_N} оказались линейными, можно сделать вывод об активационном характере исследуемого процесса. Первая зависимость имеет аррениусовский характер, вторая показывает, что коэффициент наклона первой имеет смысл.   

Была проведена экстраполяция зависимости коэффициента термической аккомодации от обратного размера кластера для температуры кластера, равной 2500~К к нулю, что соответствует бесконечно большому кластеру, то есть плоской поверхности~\cref{fig:alpha_N}. Полученное значение оказалось равным $0.21 \pm 0.06$. Именно это число имеет смысл сравнивать с экспериментальными данными, поскольку в них не определяется размер кластера, пересыщение велико и температура кластера близка к температуре плавления, потому что отвод тепла является одним из лимитирующих факторов при росте наночастиц~\cite{gurentsov_analysis_2013}. Экспериментальное значение коэффициента термической аккомодации оказалось равным 0.1. Поскольку фактическая температура и размер кластера в эксперименте невозможно оценить (хотя и есть возможность определить распределение кластеров по размерам), можно говорить о сходстве результатов с экспериментом по порядку величины. Также для температуры кластера 2500~К был проведен расчет~\cite{sipkens_examination_2015}, где было получено значение $0.23 \pm 0.03$ для плоской поверхности. С этим расчетом достигнуто сходство в пределах погрешности.

\begin{figure}[ht]
    \centerfloat{
        \includegraphics[scale=0.8]{alpha_N.png}
    \caption{Зависимости коэффициента термической аккомодации для T = 2500~K, 1000~K, 500~K. Ромб – данные расчета [8]}
    }\label{fig:alpha_N}
\end{figure}

С помощью полученных данных можно оценить коэффициент теплоотдачи в окружающую среду~\cite{grigoriev_heatmasstransfer_2005}:

\begin{align}
	k = q/\Delta T,
\end{align}

где $q$ \--- плотность теплового потока, а $\Delta T$ \--- разность температур окружающей среды и поверхности.

В этой работе рассмотрен процесс отвода теплоты от железных кластеров при помощи атомов аргона. 

\begin{align}
	k = k_B Q \nu/(T - T_a)S,
\end{align}

где $\nu = n v_T S $ \--- частота соударений атомов аргона с железными кластерами, $n$ \--- концентрация, $v_T$ \--- тепловая скорость атомов аргона, $S$ \--- площадь сечения кластера. Концентрация атомов аргона была рассчитана для нормальных условий, при которых и проводится эксперимент, тепловая скорость также соответствует 300~К.

Была проведена экстраполяция значений коэффициента теплоотдачи для конечных значений $N$ на плоскую поверхность с использованием той же зависимости от $N^{-1/3}$, что и для коэффициента термической аккоммодации. Результаты представлены в~\cref{tab:thermalconduct}.

Полученные результаты могут применяться для расчета теплоотдачи наножидкостей. Однако, процесс нагрева среды от кластеров является только одним из процессов, происходящих в наножидкостях. Вторым является передача тепла от среды к стенкам сосуда, оценка которой выходит за рамки данной работы.

\begin{table} [htbp]% Пример записи таблицы с номером, но без отображаемого наименования
    \centering
    \begin{threeparttable}% выравнивание подписи по границам таблицы
        \caption{Значения коэффициента теплоотдачи, кВт/м$^2$}%
        \label{tab:thermalconduct}%
        \begin{SingleSpace}
            \begin{tabular}{| c | c | c | c | c | c |}
                \hline
                N	& 500~К & 1000~К & 1500~К & 2000~К & 2500~К  \\ \hline
                9    & 26	& 24 	 & 23 	 & 18 	  & 18.1     \\ \hline
                27   & 34    & 26     & 24	 & 21	  & 19.7     \\ \hline
                113  & 51    & 40     & 28    & 24	  & 21 \\ \hline
                \infty & 68  & 49     & 31    & 28	  & 23.2 \\ \hline
            \end{tabular}%
        \end{SingleSpace}
    \end{threeparttable}
\end{table}
 
\section{Оценка вероятности роста кластеров}

Ещё одной задачей, рассмотренной в рамках данной работы, является расчет коэффициента прилипания атома железа к железному кластеру. Для взаимодействия между кластером и налетающим атомом использовался потенциал Леннард-Джонса с параметрами $\epsilon = 1.05*10^{-12}$ ~эрг, $\sigma = 2.95$~\AA, $r_{cut} = 5 \sigma = 14.75$~\AA. В остальном модель аналогична использованной для расчета термической аккоммодации атомов аргона. 

Коэффициент прилипания рассчитывается как доля траекторий, в результате которых атом можно было считать прилипшим. Для этого требуется выполнение 2 условий. Первое, временное, заключается в пребывании атома в зоне действия потенциала кластера более 30~пс. Это на порядок выше оценки времени, необходимого для прямого соударения атома с кластером и следующего сразу за ним вылета на расстояние, где атом и кластер не взаимодействуют. Вторым условием является энергетическое, согласно которому полная энергия атома, состоящая из кинетической энергии его движения и энергии взаимодействия с атомами кластера, должна стать отрицательной. В начальный момент времени полная энергия атома состоит только из кинетической и положительна. Атом может потерять энергию в результате соударения с атомами кластера. В случае, если полная энергия атома положительна, её достаточно, чтобы покинуть зону действия потенциала кластера.   

Далее, аналогично коэффициенту термической аккоммодации, были получены зависимости от прицельного параметра. На этих зависимостях по плато, соответствующему малым значениям прицельного параметра и наибольшему взаимодействию с кластером определяется усредненное значение коэффициента прилипания для данных температуры и размера кластера. 

Были получены зависимости коэффициента термической аккоммодации от температуры кластера. Зависимость получилась убывающей, что объясняется возрастанием энергии кластера при повышении температуры и приближением к температурам плавления и испарения.  Чем выше энергия атомов кластера, тем больше вероятность для налетающего атома получить достаточно энергии для покидания кластера. Результаты были получены для кластеров размерами 9 и 27 атомов, притом результаты для первого оказались меньше. Это следует из меньшей потенциальной энергии большего кластера. Таким образом, налетающему атому после потери части энергии при соударении нужно преодолеть больший барьер для вылета за границы дейсвтия потенциала кластера.

Была получена также и зависимость от температуры потока налетающих атомов. Чем больше начальная энергия налетающего атома, тем большую её часть необходимо потерять, чтобы прилипнуть к кластеру. Следовательно, с ростом температуры потока налетающих атомов вероятность прилипания уменьшается.

Была получена зависимость коэффициента прилипания от размера кластера. На графике по оси абсцисс отложена величина $N^{-1/3}$, поскольку поверхностная энергия возрастает как $N^{2/3}$, а объемная \--- как $N$. Полученные данные не противоречат предположению о линейности зависимости коэффициента прилипания от $N^{-1/3}$. Поскольку экспериментальные данные по величине коэффициента прилипания в вакууме не были найдены, проводилось сравнение с результатами, полученными в атмосфере различных инертных газов. Данные этой работы лучше всего совпадают с экспериментальными результатами, полученными в атмосфере ксенона, оказавшимися равными 0.1. 

\section{Выводы}

Методом молекулярной динамики рассмотерна система, состоящая из кластера и налетающего атома. Рассмотрены кластеры размером 9, 27 и 113~атомов. Поток налетающих атомов имел температуру 300~К. Получены значения количества полученной атомом энергии и коэффициента термической аккоммодации для диапазона температур кластера от 200 до 2500~К. Зависимость количества полученной атомом энергии растет с увеличением температуры кластера. Однако, это увеличение происходит по более слабому закону, чем линейный. Это объясняется уменьшением эффективности взаимодействия между атомом и кластером в силу сокращения времени взаимодействия между ними и числа соударений. 

Эту эффективность отражает коэффициент термической аккоммодации, падающий с ростом температуры. Коэффициент термической аккоммодации зависит от температуры по аррениусовскому закону. Это подтверждается как температурной зависимостью, так и оценкой этой величины из отдельных траекторий. 

Получена зависимость коэффициента термической аккоммодации от размера кластера для температуры кластера 2500~К. Результат оказался линейным от $N^{-1/3}$, что доказывает его физический смысл. Также эта экстраполяция позволила получить значение для плоской поверхности, которое находится в согласии с другим расчетом и не противоречит экспериментальным данным.

Полученные для коэффициента прилипания зависимости от размера кластера и от его температуры согласуются с экспериментом и выдвинутыми гипотезами. С ростом температуры кластера уменьшается его стабильность и уменьшается коэффициент прилипания. В случае с размером кластера наблюдается возрастание с увеличением кластера.